\documentclass[10pt,a4paper]{article}
\usepackage[T1]{fontenc}
\usepackage{lmodern,amsmath,amssymb,amsthm}
\usepackage[margin=25mm,top=19mm,bottom=21mm]{geometry}
\usepackage{xcolor}
\definecolor{linkblue}{RGB}{24,62,103}
\usepackage[colorlinks=true,linkcolor=linkblue,citecolor=linkblue,urlcolor=linkblue,pdfencoding=auto]{hyperref}
\urlstyle{same}
\newtheorem{theorem}{Theorem}
\newtheorem{lemma}[theorem]{Lemma}
\newtheorem{corollary}[theorem]{Corollary}
\newtheorem{proposition}[theorem]{Proposition}
\newcommand{\tr}{\operatorname{tr}}
\newcommand{\Id}{I}
\newcommand{\C}{\mathbb C}
\newcommand{\R}{\mathbb R}
\newcommand{\norm}[1]{\lVert #1\rVert}
\newcommand{\ket}[1]{\lvert #1\rangle}
\newcommand{\bra}[1]{\langle #1\rvert}
\newcommand{\ip}[2]{\langle #1,#2\rangle}
\newcommand{\SU}{\mathrm{SU}}
\author{Erfan Oliaei\\[2pt]\small Independent researcher}
\date{26 November 2026}
\raggedbottom
\newcommand{\ran}{\operatorname{ran}}
\newcommand{\E}{\mathbb E}
\newcommand{\ray}[1]{\lvert #1\rangle\langle #1\rvert}
\title{The exact loss threshold for two displaced on--off detectors}
\hypersetup{pdfauthor={Erfan Oliaei},pdftitle={The exact loss threshold for two displaced on--off detectors}}
\begin{document}
\maketitle
\begin{abstract}
Two displaced on--off detectors with distinct coherent centers are compatible exactly when their transmissivities sum to at most one, for every nonzero separation. The necessity follows by testing a hypothetical parent measurement on two coherent states. Equality constrains every third detector: at half transmission any three distinct collinear centers give an incompatible triple with compatible pairs. We also bound how many photon levels a fixed-radius incompatibility witness requires.
\end{abstract}
\section{Two detectors at fixed separation}
Let \(\ket a\) be the normalized single-mode coherent state of input amplitude \(a\in\C\), so \(|\langle a|b\rangle|=e^{-|a-b|^2/2}\). A detector of transmissivity \(\tau\) applies the pure-loss channel and measures the vacuum projector displaced to output amplitude \(\sqrt\tau\,a_j\). Its two outcomes are no click and click. Joint measurability means that summing effects of a single measurement produces both binary marginals.
\begin{theorem}\label{pair}
Two single-mode displaced on--off detectors of positive transmissivities
\(\tau_1,\tau_2\le1\) with distinct input coherent centers \(a_1,a_2\)
are jointly measurable exactly when \(\tau_1+\tau_2\le1\).
\end{theorem}
The no-click effect \(F_j\) has expectation
\(\langle a|F_j|a\rangle=e^{-\tau_j|a-a_j|^2}\) on coherent input \(\ket a\),
so \(F_j\ket{a_j}=\ket{a_j}\). For any parent let \(G\) be its both-no-click
effect. Each other outcome is bounded by a click effect and annihilates its detector's center state. Summing its zero cross element gives \(\langle a_1|(I-G)|a_2\rangle=0\).
On \(\ket{a_1}\) the both-no-click outcome therefore has probability \(\langle a_1|F_2|a_1\rangle\), and on \(\ket{a_2}\) it has probability \(\langle a_2|F_1|a_2\rangle\). Writing \(s=|a_1-a_2|^2>0\), Cauchy--Schwarz gives
\[
 e^{-s/2}=|\langle a_1|G|a_2\rangle|
 \le\sqrt{\langle a_1|G|a_1\rangle\langle a_2|G|a_2\rangle}
 = e^{-(\tau_1+\tau_2)s/2}.
\]
Conversely the passive splitter of Mehboudi, Rezaeinia and Rahimi-Keshari \cite[Result~1]{loss} sends fractions \(\tau_1,\tau_2\) to
two arms, discarding the rest. On each arm apply the ideal detector centered
at \(\sqrt{\tau_j}\,a_j\). This realizes the two marginals on every input.
Coincident input centers give commuting functions of the same displaced
number operator. Result~1 of \cite{loss} gives the sufficient construction for any number of arms with \(\sum_j\tau_j\le1\). Their Appendix~C obtains the equal-efficiency obstruction in the small-separation limit. Theorem~\ref{pair} gives necessity for each fixed distinct pair.

\section{What equality forces for a third detector}
Suppose \(F_1,F_2,F_3\) have a common parent and \(\tau_1+\tau_2=1\).
Cauchy--Schwarz above is an equality. In the direct sum over the underlying
parent outcomes with detectors 1 and 2 both off, the vectors
\((G_\omega^{1/2}\ket{a_1})_\omega,(G_\omega^{1/2}\ket{a_2})_\omega\)
are proportional, say \(y=cx\), because both norms are positive. If \(P\) keeps just the coordinates with detector~3 off, then \(\langle Px,Py\rangle/\langle x,y\rangle=\|Px\|^2/\|x\|^2\). All other outcomes vanish in the cross matrix element. Consequently
\[
 \frac{\langle a_1|F_3|a_2\rangle}{\langle a_1|a_2\rangle}
 = \exp[-\tau_3(\bar a_1-\bar a_3)(a_2-a_3)]\ \in[0,1].
\]
For the exponential, the loss isometry sends \(\ket a\) to
\(\ket{\sqrt{\tau_3}\,a}\otimes\ket{\sqrt{1-\tau_3}\,a}\).
Put \(b_i=\sqrt{\tau_3}\,a_i\). The environment overlap cancels from the ratio, leaving
\(\langle b_1|b_3\rangle\langle b_3|b_2\rangle/\langle b_1|b_2\rangle\).
Substitution of \(\langle u|v\rangle=\exp(\bar u v-(|u|^2+|v|^2)/2)\) gives the formula.
For three distinct collinear centers of efficiency \(1/2\), choose the middle
one as \(a_3\). The exponent is then positive, so the triple is incompatible
at every nonzero scale even though each pair is compatible. A nonreal value
of the same ratio also witnesses incompatibility for a noncollinear triple.
\section{A lower bound on the photon cutoff}
\begin{proposition}
Fix output displacement radius \(r>0\), transmissivity \(0<\tau\le1\), and
compress to photon numbers \(0,\ldots,L\). All phases at this radius have
a common parent on the compression whenever, with \(y=\tau L\),
\[
 y+2(e^{2r\sqrt y}-1)\ \le\ \min(1,e^{r^2}-1).
\]
In particular if \(0<r\le1\) and \(N\ge2\) and \(\tau=1/(N-1)\), witnessing
incompatibility by this cutoff requires \(L>r^2(N-1)/36\).
\end{proposition}
Write \(a\) for the annihilation operator, \(\hat n\) for the number operator
and \(p=e^{-r^2}\). The no-click effect centered at output \(re^{i\theta}\)
has the factorization
\[
 p\,e^{\sqrt{\tau}r e^{i\theta}a^\dagger}
       (1-\tau)^{\hat n}e^{\sqrt{\tau}r e^{-i\theta}a}.
\]
It has the required coherent-state diagonal, which determines all matrix
entries by power-series coefficients. Compression makes the expansion in
the outer exponentials finite. The term with powers \(u,v\) is bounded in
operator norm by \(p(r\sqrt y)^{u+v}/(u!v!)\), because the lowering
operator on the cutoff has norm \(\sqrt L\). The constant term from \(u=v=0\)
differs from \(pI\) by at most \(py\). Pair conjugate terms
\(u\ne v\) and write each as a cosine and a sine times Hermitian matrices.
The sum of their coefficient norms is at most twice the sum of the
original norms. The same bound covers the remaining diagonal terms. Thus
the family has the form \(pI+\sum_k t_k(\theta)A_k\) with \(|t_k|\le1\)
and \(S=\sum_k\norm{A_k}\le p[y+2(e^{2r\sqrt y}-1)]\).
For \(S>0\) select \(k\) with probability \(\norm{A_k}/S\), then measure
the effects \((I\pm A_k/\norm{A_k})/2\). On outcome \(\pm1\), the response
\(p\pm S t_k(\theta)\) gives the required effect whenever \(S\le\min(p,1-p)\).
Zero summands are omitted and \(S=0\) is immediate. For \(r\le1\),
\(y\le r^2/36\) suffices by convexity of the exponential and
\(1/36+2(e^{1/3}-1)<1\). The estimate is a necessary cutoff size for a witness; existence of a full-space witness for arbitrary fixed radius and general \(N\) is a separate question.

\paragraph{Acknowledgment.} I used OpenAI Codex to assist with literature searches, developing and checking the proofs, and preparing this manuscript.
\begin{thebibliography}{9}
\bibitem{loss} M. Mehboudi, F. Rezaeinia and S. Rahimi-Keshari, \href{https://doi.org/10.1103/s9sp-2fd7}{Measurement incompatibility under loss}, Phys. Rev. Lett. \textbf{135}, 100202 (2025). \href{https://arxiv.org/html/2411.05920v3}{Extended version, Result~1 and Appendix~C}.
\end{thebibliography}
\end{document}
